\Needspace{9\baselineskip}\begin{center}\itshape \textit{Zurückführung der Gleichungen, mit denen die Wurzeln \(\Omega\) gefunden werden, auf reine Potenzgleichungen.}\end{center}


\Needspace{6\baselineskip}\Article{359}
Die vorstehenden Untersuchungen betrafen das \textit{Auffinden} der Hilfsgleichungen; nun wollen wir eine besonders bemerkenswerte Eigenschaft ihrer \textit{Lösung} darlegen. Bekanntlich sind alle Bemühungen der größten Geometer, die allgemeine Lösung der Gleichungen über den vierten Grad hinaus oder (um genauer zu bestimmen, was ich suche) \textsc{die Zurückführung gemischter Gleichungen auf reine Potenzgleichungen} zu finden, bisher stets vergeblich gewesen, und es bleibt kaum ein Zweifel, dass dieses Problem eher etwas Unmögliches verlangt als etwas, das die Kräfte der heutigen Analysis übersteigt. (Man vergleiche unsere Bemerkungen zu diesem Gegenstand in \textit{Demonstr. nova etc. art. 9}.) Dennoch steht fest, dass es unzählige gemischte Gleichungen jedes Grades gibt, die eine solche Zurückführung auf reine Potenzgleichungen zulassen, und wir hoffen, dass es den Geometern willkommen sein wird, wenn wir zeigen, dass unsere Hilfsgleichungen immer hierzu gehören. Wegen des großen Umfangs dieser Untersuchung geben wir jedoch
\SourceLeaf{455}
hier nur die Hauptpunkte an, die zum Nachweis der Möglichkeit nötig sind, und verschieben die ausführlichere Behandlung, die dieser Gegenstand in hohem Maße verdient, auf ein anderes Mal. Vorauszuschicken sind einige allgemeine Beobachtungen über die Wurzeln der Gleichung \(x^e-1=0\), die auch den Fall einschließen, dass \(e\) zusammengesetzt ist.

I. Diese Wurzeln werden (wie aus den elementaren Lehrbüchern bekannt ist) durch \(\cos\frac{kP}e+i\sin\frac{kP}e\) dargestellt, wobei für \(k\) die \(e\) Zahlen \(0,\allowbreak 1,\allowbreak 2,\allowbreak 3\ldots e-1\) oder beliebige andere, ihnen modulo \(e\) kongruente Zahlen zu nehmen sind. Eine Wurzel wird für \(k=0\) oder allgemein für \(k\), wenn es durch \(e\) teilbar ist, \(=1\); jedem anderen Wert von \(k\) entspricht eine von \(1\) verschiedene Wurzel.

II. Da \((\cos\frac{kP}e+i\sin\frac{kP}e)^\lambda=\cos\frac{\lambda kP}e+i\sin\frac{\lambda kP}e\), ist klar: Entspricht die Wurzel \(R\) einem Wert von \(k\), der zu \(e\) teilerfremd ist, so ist in der Folge \(R,\allowbreak RR,\allowbreak R^3\), usw., das Glied an der Stelle \(e\) zwar \(=1\), alle vorhergehenden Glieder sind aber von \(1\) verschieden. Daraus folgt sofort, dass alle \(e\) Größen \(1,\allowbreak R,\allowbreak RR,\allowbreak R^3\ldots R^{e-1}\) ungleich sind; und da sie offenbar alle die Gleichung \(x^e-1=0\) erfüllen, stellen sie sämtliche Wurzeln dieser Gleichung dar.

III. Schließlich hat man unter derselben Voraussetzung \[1+R^\lambda+R^{2\lambda}\ldots+R^{\lambda(e-1)}\ =0\] für jeden ganzzahligen Wert von \(\lambda\), der nicht durch \(e\) teilbar ist; denn die Summe ist \(=\frac{1-R^{\lambda e}}{1-R^\lambda}\), und der Zähler dieses Bruches wird \(=0\), der Nenner dagegen nicht \(=0\). Ist jedoch \(\lambda\) durch \(e\) teilbar, so wird jene Summe offenbar \(=e\).

\Needspace{6\baselineskip}\Article{360}
Es sei, wie stets im Vorhergehenden, \(n\) eine Primzahl, \(g\) eine Primitivwurzel modulo \(n\) und \(n-1\) das Produkt dreier positiver ganzer Zahlen \(\alpha,\allowbreak \beta,\allowbreak \gamma\). Der Kürze wegen werden wir die Untersuchung gleich so anlegen, dass sie auch die Fälle einschließt, in denen \(\alpha\) gleich eins ist oder \(\gamma=1\) gilt; wenn \(\gamma=1\), sind anstelle der Summen \((\gamma,1),\allowbreak (\gamma,g)\), usw., die Wurzeln \([1],\allowbreak [g]\), usw., zu nehmen. Wir setzen also voraus, dass aus allen als bekannt angenommenen \(\alpha\) Summen von \(\beta\gamma\) Gliedern \((\beta\gamma,1),\allowbreak (\beta\gamma,g),\allowbreak (\beta\gamma,gg)\ldots(\beta\gamma,g^{\alpha-1})\) die Summen von \(\gamma\) Gliedern abzuleiten sind. Diese Aufgabe haben wir oben auf eine gemischte Gleichung vom Grad \(\beta\) zurückgeführt; nun wollen wir zeigen, wie man sie durch eine reine Potenzgleichung von demselben Grad löst. Zur Abkürzung bezeichnen wir die Summen \[(\gamma,1),(\gamma,g^\alpha),(\gamma,g^{2\alpha})\ldots(\gamma,g^{\alpha\beta-\alpha})\]
\SourceLeaf{456}
innerhalb von \((\beta\gamma,1)\) der Reihe nach mit \(a,\allowbreak b,\allowbreak c\ldots m\); die Summen \[(\gamma,g),(\gamma,g^{\alpha+1})\ldots(\gamma,g^{\alpha\beta-\alpha+1})\] innerhalb von \((\beta\gamma,g)\) der Reihe nach mit \(a',\allowbreak b'\ldots m'\); die Summen \[(\gamma,gg),(\gamma,g^{\alpha+2})\ldots(\gamma,g^{\alpha\beta-\alpha+2})\] der Reihe nach mit \(a'',\allowbreak b''\ldots m''\); und so weiter bis zu denjenigen, die in \((\beta\gamma,g^{\alpha-1})\) enthalten sind.

I. Nun bezeichne \(R\) unbestimmt eine Wurzel der Gleichung \(x^\beta-1=0\), und es werde angenommen, dass die Entwicklung der Potenz mit dem Exponenten \(\beta\) der Funktion \[t=a+Rb+RRc\ldots+R^{\beta-1}m\] nach den Regeln von Artikel 345 \[\begin{aligned}&N+Aa+Bb+Cc\ldots+Mm\\&{}+A'a'+B'b'+C'c'\ldots+M'm'\\&{}+A''a''+B''b''+C''c''\ldots+M''m''\\&{}+\text{usw.}=T\end{aligned}\] ergibt, wobei alle Koeffizienten \(N,\allowbreak A,\allowbreak B,\allowbreak A'\), usw., Polynome in \(R\) sein werden. Ferner werde angenommen, dass sich die Potenzen mit dem Exponenten \(\beta\) zweier anderer Funktionen \[\begin{aligned}u&=R^\beta a+Rb+RRc\ldots+R^{\beta-1}m,\\u'&=b+Rc+RRd\ldots+R^{\beta-2}m+R^{\beta-1}a\end{aligned}\] der Reihe nach in \(U\) und \(U'\) entwickeln. Aus Artikel 350 erkennt man leicht, dass, da \(u'\) aus \(t\) hervorgeht, indem man die Summen \(a,\allowbreak b,\allowbreak c\ldots m\) der Reihe nach mit \(b,\allowbreak c,\allowbreak d\ldots a\) vertauscht, \[\begin{aligned}&U'=N+Ab+Bc+Cd\ldots+Ma\\&{}+A'b'+B'c'+C'd'\ldots+M'a'\\&{}+A''b''+B''c''+C''d''\ldots+M''a''\\&{}+\text{usw.}\end{aligned}\] sein wird.

Weiter ist klar: Da \(u=Ru'\), wird \(U=R^\beta U'\) sein, sodass wegen \(R^\beta=1\) die entsprechenden Koeffizienten in \(U\) und \(U'\) gleich sind. Schließlich unterscheiden sich \(t\) und \(u\) nur dadurch, dass \(a\) in \(t\) mit der Einheit und in \(u\) mit \(R^\beta\) multipliziert wird. Daher erkennt man leicht, dass alle entsprechenden Koeffizienten (also diejenigen, mit denen dieselben Summen multipliziert werden) in \(T\) und \(U\) gleich sind und somit auch alle entsprechenden Koeffizienten in \(T\) und \(U'\). Daraus folgt endlich \(A=B=C\), usw., \(=M\);
\SourceLeaf{457}
\(A'=B'=C'\), usw., \(A''=B''=C''\), usw., usw. Dadurch lässt sich \(T\) auf eine Form wie \[N+A(\beta\gamma,1)+A'(\beta\gamma,g)+A''(\beta\gamma,gg)\ \text{usw.}\] zurückführen, wobei jeder einzelne Koeffizient \(N,\allowbreak A,\allowbreak A'\), usw., auf eine Form wie \[pR^{\beta-1}+p'R^{\beta-2}+p''R^{\beta-3}+\text{usw.}\] gebracht werden kann, in der \(p,\allowbreak p',\allowbreak p''\), usw., gegebene ganze Zahlen sind.

II. Nimmt man für \(R\) eine bestimmte Wurzel der Gleichung \(x^\beta-1=0\) (deren Lösung wir bereits als bekannt voraussetzen), und zwar eine solche, von der keine niedrigere Potenz als diejenige mit dem Exponenten \(\beta\) der Einheit gleich ist, so wird auch \(T\) eine bestimmte Größe sein, aus der sich \(t\) durch die reine Potenzgleichung \(t^\beta-T=0\) ableiten lässt. Da diese Gleichung aber \(\beta\) Wurzeln hat, nämlich \(t,\allowbreak Rt,\allowbreak RRt\ldots R^{\beta-1}t\), kann es zweifelhaft erscheinen, welche Wurzel man wählen muss. Dass diese Wahl völlig beliebig ist, wird jedoch auf folgende Weise leicht klar. Man muss sich erinnern, dass nach der Bestimmung aller Summen von \(\beta\gamma\) Gliedern die Wurzel \([1]\) nur insoweit bestimmt ist, als irgendeine der \(\beta\gamma\) in \((\beta\gamma,1)\) enthaltenen Wurzeln mit diesem Zeichen bezeichnet werden muss; und dass es daher völlig beliebig ist, welche der \(\beta\) Summen, aus denen \((\beta\gamma,1)\) besteht, wir mit \(a\) bezeichnen wollen. Nehmen wir nun an, dass bei Bezeichnung einer bestimmten Summe mit \(a\) \(t=\mathfrak{T}\) wird. Dann erkennt man leicht: Will man nachher die eben mit \(b\) bezeichnete Summe mit \(a\) bezeichnen, so werden diejenigen, die vorher \(c,\allowbreak d\ldots a,\allowbreak b\) waren, nun \(b,\allowbreak c\ldots m,\allowbreak a\), und der Wert von \(t\) wird jetzt \(=\frac{\mathfrak{T}}R=\mathfrak{T}R^{\beta-1}\). Ebenso wird, wenn man mit \(a\) die ursprünglich mit \(c\) bezeichnete Summe bezeichnen will, der Wert von \(t\) \(\mathfrak{T}R^{\beta-2}\), und so weiter: \(t\) kann als gleich jeder der Größen \(\mathfrak{T},\allowbreak \mathfrak{T}R^{\beta-1},\allowbreak \mathfrak{T}R^{\beta-2}\), usw., \(\text{d.\,h.}\) jeder Wurzel der Gleichung \(x^\beta-T=0\) angesehen werden, je nachdem man annimmt, dass diese oder jene in \((\beta\gamma,1)\) enthaltene Summe mit \((\gamma,1)\) bezeichnet wird. \(Q.\ E.\ D.\)

III. Nachdem die Größe \(t\) auf diese Weise bestimmt ist, muss man \(\beta-1\) andere suchen, die aus \(t\) hervorgehen, wenn man in seinem Ausdruck für \(R\) nacheinander \(RR,\allowbreak R^3,\allowbreak R^4\ldots R^\beta\) einsetzt, nämlich \[\begin{aligned}t'&=a+RRb+R^4c\ldots+R^{2\beta-2}m,\\t''&=a+R^3b+R^6c\ldots+R^{3\beta-3}m\ \text{usw.}\end{aligned}\] Die letzte ist bereits bekannt, da sie offenbar \(=a+b+c\ldots+m=(\beta\gamma,1)\) wird; die übrigen lassen sich folgendermaßen finden. Entwickelt man nach den Regeln von Artikel 345, ebenso wie vorher in I \(t^\beta\), das Produkt \(t^{\beta-2}t'\), so wird man mit einer dem vorigen Verfahren völlig entsprechenden Methode beweisen, dass sich das Ergebnis auf eine Form wie
\SourceLeaf{458}
{\postdisplaypenalty=10000 \[\mathfrak{N}+\mathfrak{A}(\beta\gamma,1)+\mathfrak{A}'(\beta\gamma,g)+\mathfrak{A}''(\beta\gamma,gg)\ \text{usw.}=T'\]} zurückführen lässt, in der \(\mathfrak{N},\allowbreak \mathfrak{A},\allowbreak \mathfrak{A}'\), usw., Polynome in \(R\) sind, sodass \(T'\) eine bekannte Größe ist, aus der man \(t'=\frac{T'tt}{T}\) erhält. Nimmt man in genau derselben Weise an, dass die Entwicklung des Produkts \(t^{\beta-3}t''\) \(T''\) ergibt, so wird dieser Ausdruck eine ähnliche Form haben, und aus seinem bekannten Wert wird man daher \(t''\) durch die Gleichung \(t''=\frac{T''t^3}{T}\) ableiten; ebenso wird \(t'''\) durch eine Gleichung wie \(t'''=\frac{T''''t^4}{T}\) gefunden, in der \(T'''\) eine bekannte Größe ist, und so weiter.

Diese Methode wäre nicht anwendbar, wenn \(t=0\) eintreten könnte, woraus auch \(T=T'=T''\), usw., \(=0\) sein müsste; doch lässt sich beweisen, dass dies unmöglich ist, auch wenn wir den Beweis wegen seiner Länge hier weglassen müssen. — Es gibt auch besondere Kunstgriffe, mit denen sich die Brüche \(\frac{T'}T,\allowbreak \frac{T''}T\), usw., in \textit{Polynome} in \(R\) verwandeln lassen, sowie kürzere Methoden, im Fall \(\alpha=1\) die Werte von \(t',\allowbreak t''\), usw., zu finden. Dies alles müssen wir hier ohne weitere Erörterung übergehen.

IV. Sobald schließlich \(t,\allowbreak t',\allowbreak t''\), usw., gefunden sind, hat man nach Beobachtung III des vorigen Artikels sofort \(t+t'+t''+\text{usw.}=\beta a\), womit der Wert von \(a\) bekannt ist; daraus lassen sich nach Artikel 346 die Werte aller übrigen Summen von \(\gamma\) Gliedern ableiten. — Die Werte von \(b,\allowbreak c,\allowbreak d\), usw., können auch aus den folgenden Gleichungen gewonnen werden, deren Grund jedem aufmerksamen Leser leicht ersichtlich sein wird: \[\begin{aligned}\beta b&=R^{\beta-1}t+R^{\beta-2}t'+R^{\beta-3}t''+\text{usw.}\\\beta c&=R^{2\beta-2}t+R^{2\beta-4}t'+R^{2\beta-6}t''+\text{usw.}\\\beta d&=R^{3\beta-3}t+R^{3\beta-6}t'+R^{3\beta-9}t''+\text{usw. usw.}\end{aligned}\]

Von der großen Zahl der Beobachtungen zur vorstehenden Untersuchung berühren wir hier nur eine. Was die Lösung der reinen Potenzgleichung \(x^\beta-T=0\) angeht, so hat \(T\) offenbar in den meisten Fällen einen imaginären Wert \(P+iQ\), sodass jene Lösung bekanntlich teils von der Teilung eines Winkels (dessen Tangens \(=\frac QP\) ist), teils von der Teilung eines Verhältnisses (der Einheit zu \(\surd(PP+QQ)\)) in \(\beta\) Teile abhängt. Dabei ist sehr bemerkenswert (was wir hier jedoch nicht weiter ausführen), dass sich der Wert von \({}^{\beta}\!\surd(PP+QQ)\) stets \textit{rational} durch schon bekannte Größen ausdrücken lässt, sodass außer der Ausziehung einer Quadratwurzel zur Lösung \textit{nur} die Teilung eines Winkels erforderlich ist, \(\text{z.\,B.}\) für \(\beta=3\) nur die Dreiteilung eines Winkels.

Da schließlich nichts hindert, \(\alpha=1,\allowbreak \gamma=1\) und somit
\SourceLeaf{459}
\(\beta=n-1\) zu setzen, ist klar, dass sich die Lösung der Gleichung \(x^n-1=0\) sogleich auf die Lösung der reinen Potenzgleichung vom Grad \(n-1\), nämlich \(x^{n-1}-T=0\), zurückführen lässt, wobei \(T\) durch die Wurzeln der Gleichung \(x^{n-1}-1=0\) bestimmt wird. Daraus folgt mithilfe der eben gemachten Beobachtung, dass die Teilung des ganzen Kreises in \(n\) Teile \(1^\circ\) die Teilung des ganzen Kreises in \(n-1\) Teile, \(2^\circ\) die Teilung eines anderen Bogens, der nach Ausführung jener Teilung konstruiert werden kann, in \(n-1\) Teile und \(3^\circ\) die Ausziehung einer einzigen Quadratwurzel erfordert; und es lässt sich tatsächlich zeigen, dass diese stets \(\surd n\) ist.
